\scanpage{28}
\[
\operatorname{Coker}T_{\deleted{\unsure{p}}p}(\widetilde\varphi)_{\mathrm{\acute et}}\simeq
(\operatorname{Ker}\deleted{T}\widetilde\varphi)(p)_{\mathrm{\acute et}}
\]
Or on sait que $\operatorname{Ker}\widetilde\varphi$ est autodual, donc
$(\operatorname{Ker}\widetilde\varphi)(p)_{\mathrm{\acute et}}$ est dual de $(\operatorname{Ker}\widetilde\varphi(p))_{\mathrm{rad}}$, donc
\unsure{avec|de|comme} $=$ cardinal, d'où la relation
\[
\begin{aligned}
\bigl(\operatorname{card}(\operatorname{Ker}\widetilde\varphi)(p)_{\mathrm{\acute et}}\bigr)^2
&=\operatorname{card}\operatorname{Ker}(\widetilde\varphi(p))\\
&\hspace{2em}\shortparallel\\[-1ex]
&\quad(\deg\widetilde\varphi)(p).
\end{aligned}
\]
d'où la relation envisagée \ldots\ldots\ Cela définit
bien un isom canonique
\[
\boxed{\det T_p(A)_{\mathrm{\acute et}}\xrightarrow[\alpha_A]{\sim}\det T_p(A^*)_{\mathrm{\acute et}}}\tag*{(*)}
\]
compatible avec \unsure{changmt|changt} de base, et isom. en $A$.

[Si $u$ est un autom de $A$, et $u^\vee=u^{*\, -1}$ son
contragrédient, on aura $\alpha\,\det T_p(u)_{\mathrm{\acute et}}=\operatorname{d\acute et}T_p(u^\vee)_{\mathrm{\acute et}}\alpha$,
donc en identifiant \unsure{des|les} dét's \unsure{des|à des} élts de
$\mathbb Z_p$ \unsure{$\times$|ici}, on aura
$\operatorname{d\acute et}T_p(u)_{\mathrm{\acute et}}=\underset{\substack{\shortparallel\\(\operatorname{d\acute et}T_p(u^*))^{-1}}}{\operatorname{det}T_p(u^{*\, -1})_{\mathrm{\acute et}}}$
donc
\[
\boxed{\det T_p(u)_{\mathrm{\acute et}}\inserted{\det}T_p(u^*)_{\mathrm{\acute et}}=1}
\qquad\text{\uline{$u$ un autom de $A$}} \quad]
\]

\margin{\unsure{2|Z}}
Attention, si $u$ est un endom de $A$, on aura
en général $\det T_p(u)\ne\det T_p(u^*)$
% archive sheet number: 27
